\begin{abstract}
Visual counterfactual explanation methods for medical image classifiers are conventionally scored by Counterfactual Validity (CV): the rate at which a generated image flips the classifier's decision. We make two contributions that pull in opposite directions, and argue the second is the more consequential. First, we present PPO-CF, a system in which a PPO agent schedules the sampling procedure of a frozen, LoRA-adapted latent diffusion generator (selecting per-step noise strength, classifier- and preservation-guidance scales, refinement-pass count, guidance ramp-up point and stopping time) as an alternative to one hand-tuned schedule for every image, which may fit a heterogeneous population poorly; editing is largely confined to a dilated lung region by a soft, RePaint-style background lock. Second, we show that CV, computed as it usually is with the same classifier that guided generation, cannot by itself establish whether the edit is class-relevant. We introduce two measurements that do not share that circularity: a Lesion Hit Rate (LHR) anchored to radiologist-drawn opacity boxes, and a held-out transfer ratio that re-scores counterfactuals with classifiers that never participated in guidance. Under these, the picture changes substantially. On the RSNA Pneumonia Detection Challenge, our committed configuration attains a mean CV of 0.838 and a 2.67$\times$ LHR lift over the annotation-area null ($p = 4.2 \times 10^{-53}$), yet on the most independent validator held-out transfer for n$\to$p remains in the adversarial regime (0.046 n$\to$p, 0.186 p$\to$n, against a 0.15 threshold). A five-arm study of trained policies (one training run per arm) produces two negative results: an annealed sparsity reward fails at its own objective and costs lesion targeting ($d = 0.512$), and a Grad-CAM spatial constraint collapses the base task entirely ($d = 1.205$, CV $0.838 \to 0.017$). The committed ensemble-guided policy also targets lesions better than an earlier single-classifier policy (paired $d = 0.482$), but that policy predates our split and seeding controls (\S\ref{sec:controls}) and likely differs in other training settings as well, so the gain cannot be attributed to ensembling alone.
% TODO: "likely differs" rests on development history (reward weights, a1 range, anchoring), unverified for the trained arm
We argue that reporting CV without a ground-truth- or independence-anchored companion metric can substantially overstate progress, as it does for our own system.
\end{abstract}
